% Part: incompleteness
% Chapter: incompleteness-provability
% Section: lob-thm

\documentclass[../../../include/open-logic-section]{subfiles}

\begin{document}

\olfileid[tr]{inc}{inp}{lob}

\olsection{L\"ob Teoremi}

Bir kuram $\Th{T}$ verilsin. Bu kuramın G\"odel tümcesi,
$\lnot \OProv[\Th{T}](y)$'nin bir sabit noktasıdır; yani
\[
\Th{T} \Proves \lnot \OProv[\Th{T}](\gn{!G}) \liff !G.
\]
olacak biçimde !!a{sentence}~$!G$'dir. Bu tümce !!{derivable} değildir;
çünkü $\Th{T} \Proves !G$ olsaydı (a) (1) numaralı !!{derivability} koşulu
uyarınca $\Th{T} \Proves \OProv[\Th{T}](\gn{!G})$ olurdu; (b)
$\Th{T} \Proves !G$ ile
$\Th{T} \Proves \lnot \OProv[\Th{T}](\gn{!G}) \liff !G$ birlikte
$\Th{T} \Proves \lnot \OProv[\Th{T}](\gn{!G})$ sonucunu verirdi ve böylece
$\Th{T}$ tutarsız olurdu. Şimdi $\OProv[\Th{T}](y)$'nin sabit noktasının,
yani
\[
\Th{T} \Proves \OProv[\Th{T}](\gn{!H}) \liff !H.
\]
olacak biçimde !!a{sentence}~$!H$'nin durumunu sormak doğaldır. $!H$
!!{derivable} olsaydı (1) koşulu uyarınca
$\Th{T} \Proves \OProv[\Th{T}](\gn{!H})$ olurdu; fakat yukarıdaki denkliğe
modus ponens uygularsak da aynı sonucu elde ederiz. Dolayısıyla, en azından
G\"odel tümcesi durumundaki argümanla, $\Th{T}$'nin tutarsız olduğu sonucunu
elde etmeyiz. Elbette bu, $\Th{T}$'nin $!H$'yi \emph{gerçekten}
!!{derive} ettiğini göstermez.

Soruyu biraz genelleştirirsek ilerleme kaydedebiliriz. Sabit nokta denkliğinin
soldan sağa yönü, $\OProv[\Th{T}](\gn{!H}) \lif !H$, \emph{yansıma ilkesi}
denen genel bir şemanın örneğidir:
$\OProv[\Th{T}](\gn{!A}) \lif !A$. Bu adın verilmesinin nedeni bir anlamda
$\Th{T}$'nin !!{derive} edebildikleri üzerine ``düşünebildiğini'' ifade
etmesidir; ilke özünde, her $!A$ için ``$\Th{T}$, $!A$'yı !!{derive}
edebiliyorsa $!A$ doğrudur'' der. Bu elbette yalnızca sağlam kuramlar için
doğrudur ve kuramların genel olarak ilkenin her örneğini !!{derive}
etmeyeceğini düşündürür. Öyleyse (yeterince güçlü ve !!{derivability}
koşullarını sağlayan) bir kuram hangi örnekleri !!{derive} edebilir? Kuşkusuz
$!A$'nın kendisinin !!{derivable} olduğu bütün örnekleri. Aşağıdaki sonucun
gösterdiği gibi yalnızca bunları.

\begin{thm}\ollabel{thm:lob}
$\Th{T}$, $\Th{Q}$'yu genişleten !!a{axiomatizable} kuram olsun ve
$\OProv[\Th{T}](y)$'nin
% [tr source emendation] The P1--P3 conditions are introduced in the `prc`
% section; upstream points to `2in`, the following second-theorem section.
\olref[prc]{sec}'deki P1--P3 koşullarını sağlayan bir formül olduğunu
varsayın. $\Th{T}$, $\OProv[\Th{T}](\gn{!A}) \lif !A$ formülünü !!{derive}
ederse gerçekten de $!A$'yı !!{derive} eder.
\end{thm}

Başka türlü söylersek $\Th{T} \Proves/ !A$ ise
$\Th{T} \Proves/ \OProv[\Th{T}](\gn{!A}) \lif !A$ olur. Bu sonuç L\"ob
teoremi olarak bilinir.

\begin{explain}
L\"ob teoreminin ispatı için yol gösterici fikir, Noel Baba'nın var olduğunu
gösteren zekice bir ispattır. (Bu sonucu beğenmiyorsanız istediğiniz başka bir
sonucu koymakta özgürsünüz.) İspat şöyledir:
\begin{enumerate}
\item $X$, ``$X$ doğruysa Noel Baba vardır'' tümcesi olsun.
\item $X$'in doğru olduğunu varsayın.
\item O zaman söylediği şey geçerlidir; yani $X$ doğruysa Noel Baba vardır.
\item $X$'in doğru olduğunu varsaydığımızdan (2) ve (3)'e modus ponens
  uygulayarak Noel Baba'nın var olduğu sonucuna ulaşabiliriz.
\item (4)'teki ``Noel Baba vardır'' sonucunu, (2)'deki ``$X$ doğrudur''
  varsayımından türetmeyi başardık. Koşullu ispatla ``$X$ doğruysa Noel Baba
  vardır'' önermesini göstermiş olduk.
\item Fakat bu tam olarak $X$ tümcesidir. Demek ki $X$'in doğru olduğunu
  gösterdik.
\item O zaman yukarıdaki (2)--(4) argümanına göre Noel Baba vardır.
\end{enumerate}
Bu fikri, ``doğrudur'' yerine ``!!{derivable}{}dir'', ``Noel Baba vardır''
yerine de $!A$ koyarak biçimselleştirmek L\"ob teoreminin ispatını verir.
Hile, sabit nokta lemmasını $\OProv[\Th{T}](y) \lif !A$ !!{formula}'na
uygulamaktır. Bunun sabit noktası, yukarıdaki taslakta yer alan $X$ tümcesine
karşılık gelir.
\end{explain}

\begin{proof}[\olref{thm:lob} İspatı]
$!A$, !!a{sentence} olsun ve $\Th{T}$,
$\OProv[\Th{T}](\gn{!A}) \lif !A$ formülünü !!{derive} etsin. $!B(y)$,
$\OProv[\Th{T}](y) \lif !A$ !!{formula} olsun. Sabit nokta lemmasıyla
!!a{sentence}~$!D$ bulunsun ve $\Th{T}$, $!D \liff !B(\gn{!D})$
formülünü !!{derive} etsin.
O zaman aşağıdakilerin her biri $\Th{T}$ içinde !!{derivable}{}dir:
\begin{align}
  & !D \liff (\OProv[\Th{T}](\gn{!D}) \lif !A) \ollabel{L-1}\\
  & \qquad \text{$!D$, $!B(y)$'nin bir sabit noktasıdır}\notag \\
  & !D \lif (\OProv[\Th{T}](\gn{!D}) \lif !A) \ollabel{L-2}\\
  & \qquad\text{\olref{L-1}'den}\notag\\
  & \OProv[\Th{T}](\gn{!D \lif (\OProv[\Th{T}](\gn{!D}) \lif !A)}) \ollabel{L-3}\\
  & \qquad \text{\olref{L-2}'den P1 koşulu aracılığıyla}\notag \\
  & \OProv[\Th{T}](\gn{!D}) \lif \OProv[\Th{T}](\gn{\OProv[\Th{T}](\gn{!D}) \lif !A})
  \ollabel{L-4}\\
  &\qquad \text{\olref{L-3}'ten P2 koşulu kullanılarak}\notag \\
  & \OProv[\Th{T}](\gn{!D}) \lif (\OProv[\Th{T}](\gn{\OProv[\Th{T}](\gn{!D})}) \lif \OProv[\Th{T}](\gn{!A})) \ollabel{L-5}\\
  &\qquad \text{\olref{L-4}'ten P2 yeniden kullanılarak} \notag\\
& \OProv[\Th{T}](\gn{!D}) \lif \OProv[\Th{T}](\gn{\OProv[\Th{T}](\gn{!D})}) \ollabel{L-6}\\
  & \qquad\text{P3 !!{derivability} koşulu aracılığıyla} \notag\\
  & \OProv[\Th{T}](\gn{!D}) \lif \OProv[\Th{T}](\gn{!A}) \ollabel{L-7} \\
  &\qquad\text{\olref{L-5} ve \olref{L-6}'dan}\notag\\
  & \OProv[\Th{T}](\gn{!A}) \lif !A \ollabel{L-8}\\
  &\qquad\text{teoremin varsayımından} \notag\\
  & \OProv[\Th{T}](\gn{!D}) \lif !A \ollabel{L-9}\\
  &\qquad\text{\olref{L-7} ve \olref{L-8}'den}\notag\\
  & (\OProv[\Th{T}](\gn{!D}) \lif !A) \lif !D \ollabel{L-10}\\
  & \qquad \text{\olref{L-1}'den}\notag \\
  & !D \ollabel{L-11}\\
  & \qquad\text{\olref{L-9} ve \olref{L-10}'dan}\notag \\
  & \OProv[\Th{T}](\gn{!D}) \ollabel{L-12}\\
  & \qquad\text{\olref{L-11}'den P1 koşulu aracılığıyla}\notag \\
  & !A \qquad\qquad\text{\olref{L-8} ve \olref{L-12}'den}\notag
\end{align}
\end{proof}

L\"ob teoremi eldeyken, (P1)--(P3) koşullarını sağlayan !!a{derivability}
yüklemi bulunan kuramlar için ikinci eksiklik teoreminin kısa bir ispatı
vardır: $\Th{T} \Proves \OProv[\Th{T}](\gn{\lfalse}) \lif \lfalse$ ise
$\Th{T} \Proves \lfalse$ olur. $\Th{T}$ tutarlıysa
$\Th{T} \Proves/ \lfalse$ olur. Öyleyse
$\Th{T} \Proves/ \OProv[\Th{T}](\gn{\lfalse}) \lif \lfalse$, yani
$\Th{T} \Proves/ \OCon[\Th{T}]$ olur. Teoremi ayrıca
$\OProv[\Th{T}](x)$'in sabit noktası $!H$'nin !!{derivable} olduğunu göstermek
için de uygulayabiliriz. Çünkü
\begin{align*}
  \Th{T} & \Proves \OProv[\Th{T}](\gn{!H}) \liff !H\\
  \intertext{ve özellikle}
    \Th{T} & \Proves \OProv[\Th{T}](\gn{!H}) \lif !H
\end{align*}
olur; dolayısıyla L\"ob teoremi uyarınca $\Th{T} \Proves !H$ olur.

% Going in the other direction, for homework I
% may ask you to work through a short proof of L\"ob's theorem, using
% the second incompleteness theorem instead of the fixed-point lemma.

\begin{prob}
$\Th{T}$ hesaplanabilir biçimde aksiyomlaştırılmış bir kuram ve
$\OProv[\Th{T}]$, $\Th{T}$ için !!a{derivability} yüklemi olsun. Aşağıdaki
dört önermeyi ele alın:
% [tr source emendation] Upstream drops the \Th{} theory macro throughout
% this list; the surrounding notation consistently denotes the theory \Th{T}.
\begin{enumerate}
\item $\Th{T} \Proves !A$ ise $\Th{T} \Proves \OProv[\Th{T}](\gn{!A})$ olur.
\item $\Th{T} \Proves !A \lif \OProv[\Th{T}](\gn{!A})$.
\item $\Th{T} \Proves \OProv[\Th{T}](\gn{!A})$ ise $\Th{T} \Proves !A$ olur.
\item $\Th{T} \Proves \OProv[\Th{T}](\gn{!A}) \lif !A$
\end{enumerate}
Bu önermelerin her biri hangi koşullar altında doğrudur?
\end{prob}

\end{document}
